\section{Introduction}
A numerical command experiment asks a finite register to retain enough
information to answer one selected observable after a word of known
length. The number of labels at a single cut is a static positive
factorization problem. Using small factorizations at every cut requires
one compatible chain of stochastic updates, and is a different problem.
For orthogonal command dynamics, orbit geometry and a conditional-mean
entropy budget quantify this difference.

The hypothesis should be a property of the experiment. A presentation
by a group with an unused compact factor can change the full regular
representation gap without changing a single response probability.
Likewise, decomposing an isotypic space into equivalent irreducible
copies can change individual seed projections without changing the
experiment. We remove both choices: expansion is measured in the
actual action, and activation is optimized over isometric intertwiners.
The latter optimization has an exact commutant covariance formula for
its unweighted magnitude. These changes also permit nonspanning seeds
and experiments with nonexpanding, genuinely acting factors.

Fix orthogonal matrices $A_a$ on $V=\R^D$, indexed by a finite alphabet,
and a nonempty finite set $\mathcal X$ of unit seeds. For a word
$w=a_1\cdots a_N$ let $A_w=A_{a_N}\cdots A_{a_1}$. With a selected
coordinate $j$ and a binary answer $B$, the law is
\begin{equation}\label{eq:intro-law41}
 \Pp(B=b\mid x,w,j)=\frac{1+b\rho\langle e_j,A_wx\rangle}{2}.
\end{equation}
The error is the maximum binary total-variation error over all seeds,
words and queries. The width $W_{N,\varepsilon}$ is the minimum of the
maximum number of available persistent labels in a realizing machine,
including the final two-label answer cut. The horizon and epoch,
nonuniform row tables, arithmetic and exact atomic sampling are free.
A different machine may be used for every $N$. These are finite positive
realizations, not uniform space-bounded algorithms.

Let $H=\overline{\langle A_a\rangle}$ be connected and put
$R=\operatorname{span}(H\mathcal X)$, $r=\dim R$. Its canonical isotypic
decomposition has invariant part $R_0$ and nontrivial types $\lambda$.
For one irreducible model $E_\lambda$ of each active type, define
\[
 s_\lambda=\min_{\|v\|=1}\dim(Hv),\qquad
 \sigma_{\mathcal X}=\max_\lambda s_\lambda.
\]
The maximum is zero if there is no nontrivial type. For an isometric
intertwiner $J:E_\lambda\to R_\lambda$, let
$B(J)=\|J^*\|_{\ell^\infty_D\to\ell^2}$, and define
\[
 c_\lambda=\max_{x\in\mathcal X,\ J}
                     \frac{\|J^*x\|}{B(J)}.
\]
All these quantities belong to the represented experiment and its
specified query geometry. For a command probability $p$, the action
gap on $E_\lambda$ is
$g_\lambda=1-\|P_\lambda-\Pi_\lambda\|^2$, where $P_\lambda$
averages the commands on sphere functions and $\Pi_\lambda$ is Haar
projection onto the invariant functions.

\begin{theorem}[Intrinsic dominant-action law]\label{thm:main}
Suppose $0<\rho\le\min\{1/10,1/(8r)\}$. If an active type
$\lambda_*$ attaining $\sigma_{\mathcal X}$ has $g_{\lambda_*}>0$
and $0\le\varepsilon<\rho c_{\lambda_*}/2$, then
\begin{equation}\label{eq:main41}
                 W_{N,\varepsilon}=\Theta(N^{\sigma_{\mathcal X}/2}).
\end{equation}
No action gap is required on the other types. The lower bound permits
arbitrary hidden states and time-dependent stochastic rows. An exact
upper realization uses common command rows and a horizon-dependent
terminal decoder, with its randomized branch identity charged to the
register. If all active vectors are invariant, constant width suffices.
\end{theorem}

The action gap is identical on sphere functions and on Lebesgue-square-
integrable functions on the representation space. It has an exact
harmonic-block norm description. A full regular-representation gap of
the effective group remains a sufficient certificate, but is no longer
required by the headline. Proposition~\ref{prop:mixed41} makes the
extension strict: ten rational orthogonal commands act faithfully on
$\R^3\oplus\R^2$, generate $\SO(3)\times\SO(2)$, and have no full
group norm gap, while their numerical width is $\Theta(N)$. The
three-dimensional constituent expands; the circle constituent does not.

We also compute rather than merely name the orbit invariant.
Theorem~\ref{thm:adjoint41} determines every compact simple adjoint
minimum by a single-node root deletion, including all exceptional
types and the minimizing centralizers. Theorem~\ref{thm:exterior41}
determines the invariant throughout the second exterior-power
fundamental family. In its $\SU(6)$ example the minimizing orbit is
$\SU(6)/\operatorname{Sp}(3)$ of dimension fourteen; the decomposable
highest-weight orbit has dimension seventeen. This illustrates why
projective highest-weight geometry cannot be substituted for the
real-vector orbit required by the memory theorem.

\begin{center}
\begin{tabular}{p{0.18\textwidth}p{0.73\textwidth}}
\toprule
Notation & Meaning\\
\midrule
$H,R$ & Effective compact matrix group and invariant span of the seeds.\\
$R_\lambda,E_\lambda$ & Canonical isotypic part and one irreducible model.\\
$s_\lambda,\sigma_{\mathcal X}$ & Minimum orbit dimension and dominant active minimum.\\
$g_\lambda,c_\lambda$ & Intrinsic action gap and query-calibrated activation.\\
$K_t,W_{N,\varepsilon}$ & Available labels at a cut and optimal whole-run peak.\\
\bottomrule
\end{tabular}
\end{center}

The analytic foundation is the orbital entropy occupation theorem,
whose complete proof is retained. Uniform orbit charts and contracted
orbit hulls supply its geometric lower and constructive upper partners.
Those results are inherited; the present extensions are the intrinsic
action formulation, canonical activation on arbitrary seed spans,
dominant-only expansion, and the resulting faithful nongapped example.
Root centralizers, skew congruence normal forms, covering geometry and
positive attenuating embeddings are classical tools, not new names
for discoveries. The original projective and arithmetic applications
retain their separate signal ranges and external assumptions in the
appendices. No main result relies on the optional numerical LPS input.
